\documentclass[12pt]{article}
\usepackage[T1]{fontenc}
\usepackage[utf8]{inputenc}
\usepackage{lmodern}
\usepackage{textcomp}
\usepackage{enumitem}
\usepackage{geometry}
\usepackage{graphicx}
\usepackage{amsmath, amssymb}
\geometry{margin=1in}
\begin{document}
\begin{center}
Constitutional Law - Undergraduate Level Assessment 23 \
Total Marks: 40 \
Answer all questions.
\end{center}

\section*{Multiple Choice (10 marks)}
\begin{enumerate}[label=\arabic*.]
\item Which of the following best describes Dworkin's distinction between principles and policies?
\begin{enumerate}[label=(\alph*)]
    \item Principles are legislative, policies are precedents.
    \item Principles describe rights, policies describe duties.
    \item Principles are democratic, policies are autocratic.
    \item Principles describe rights, policies describe goals.
\end{enumerate}
\item What is 'unilateral acts'?
\begin{enumerate}[label=(\alph*)]
    \item They are acts that States perform as practice in the context of custom
    \item They are acts creating unilateral legal obligations to the acting State
    \item Unilateral acts are simply political acts of State devoid of any legal effect
    \item Unilateral acts are those that State perform in order to be bound by a treaty
\end{enumerate}
\item Functions of the law include all but which of the following?
\begin{enumerate}[label=(\alph*)]
    \item maximizing individual freedom
    \item providing a basis for compromise
    \item keeping the peace
    \item promoting the principles of the free enterprise system
\end{enumerate}
\item Is piracy under international (jure gentium) law subject to universal jurisdiction?
\begin{enumerate}[label=(\alph*)]
    \item Piracy jure gentium is subject to flag State jurisdiction
    \item Piracy jure gentium is subject to universal jurisdiction
    \item Piracy jure gentium is subject to port State jurisdiction
    \item Piracy jure gentium is subject to nationality-based jurisdiction
\end{enumerate}
\item Which are the formal sources of international law?
\begin{enumerate}[label=(\alph*)]
    \item Custom, treaties and judicial decisions
    \item Custom, general principles of law and theory
    \item Treaties, custom and general principles of law
    \item Treaties, custom and General Assembly Resolutions
\end{enumerate}
\end{enumerate}

\section*{True / False (10 marks)}
\textit{Answer True or False}
\begin{enumerate}[label=\arabic*.]
\setcounter{enumi}{5}
\item True or False: The Sociological School of jurisprudence asserts that law serves to achieve and advance specific social goals.
\item True or False: The prohibition of torture was the first rule of jus cogens explicitly recognized by the International Court of Justice.
\item True or False: There is no relationship between law and morality, as they operate independently of each other in society.
\item True or False: Kelsen's pure theory of law emphasizes the separation of law from religion, ethics, sociology, and history.
\item True or False: The Analytical School of jurisprudence asserts that law is purely a system of logical rules independent of social influences.
\end{enumerate}

\section*{Long Form Response (20 marks)}
\textit{Answer each question with a clear and complete explanation.}
\begin{enumerate}[label=\arabic*.]
\setcounter{enumi}{10}
\item Discuss the role of consent as a circumstance that can preclude the wrongfulness of state conduct in constitutional law, highlighting the conditions for valid consent.
\item Discuss the Marxist critique of individual rights within the context of constitutional law, focusing on their relationship to capitalism and their role in a socialist society.
\end{enumerate}

\vfill
\noindent\textit{End of Paper}
\end{document}